\documentclass[journal]{IEEEtran}

\usepackage{amsmath,amssymb}
\usepackage{graphicx}
\usepackage{booktabs}
\usepackage{multirow}
\usepackage{url}
\usepackage[hidelinks]{hyperref}

\title{HyperLIF Learning: Sparse Temporal State in Continuous-Reset Spiking Networks}

\author{Gandhi Alexis Sinhue Contreras Torres, Steven Sandoval, and Luis Rodolfo Garcia Carrillo%
\thanks{Proposal scaffold prepared on 2026-07-19. Replace author metadata and acknowledgments before submission.}}

\begin{document}
\maketitle

\begin{abstract}
Spiking neural networks promise event-driven temporal computation, but standard leaky integrate-and-fire neurons expose only discontinuous threshold-reset events and often require surrogate gradients to train. HyperLIF augments the LIF membrane with a bounded continuous reset gate that preserves spike/reset semantics while exposing an internal state variable. This paper proposes a controlled cross-architecture study of HyperLIF and Adaptive HyperLIF in temporal learning. Across feedforward, convolutional, recurrent, delay, and residual spiking architectures, we evaluate whether the gate state improves temporal-task accuracy, readout interpretability, and spike efficiency under matched data, training, and seed budgets. The current evidence supports a disciplined claim: HyperLIF is not a generic leaderboard mechanism, but a sparse temporal-state mechanism whose gate is decodable, whose advantage grows with temporal content, and whose efficiency/readout trade can be tuned across architectures and frameworks.
\end{abstract}

\begin{IEEEkeywords}
Spiking neural networks, HyperLIF, Adaptive HyperLIF, surrogate gradients, temporal learning, event-based vision, sparse neural computation.
\end{IEEEkeywords}

\section{Introduction}
The rejected TNNLS manuscript made the theory too central before the practical value was clear. This paper inverts the order. The reader first sees the computational question: when a spiking neuron exposes a continuous gate state $r(t)$, does that state improve temporal learning or sparse readout under controlled budgets?

The central thesis is that HyperLIF contributes an interpretable temporal state absent in LIF. This state is useful as a readout, as an efficiency dial, and as a bridge to local learning and control. The contribution is empirical and mechanistic, not a claim that exact differentiability eliminates surrogate training.

\section{Related Work and Positioning}
This section must explicitly answer the TNNLS novelty critique.

Required comparison table:
\begin{itemize}
  \item Smooth gated continuous RNNs: coRNN, UnICORNN, LEM.
  \item Spiking neuron baselines: LIF, ALIF/LSNN, PLIF, GLIF or learnable-threshold LIF.
  \item Continuous/hybrid spiking models: Izhikevich, spike response models, adaptive-threshold neurons.
  \item Deep SNN training: surrogate-gradient LIF, SEW-ResNet, delay SNN, Spikformer if feasible.
\end{itemize}

The distinction to defend: HyperLIF is a spike/reset neuron with measurable event rate, readout sparsity, and a continuous gate variable, not merely a smooth RNN cell.

\section{Model}
Present only the equations needed to understand the experiments:
\begin{align}
\tau_m \dot{v} &= -v + J - G(r, J)(v-v_{\mathrm{reset}}),\\
\tau_r \dot{r} &= \sigma(\alpha(v-(v_{\mathrm{th}}-\gamma r))) - r.
\end{align}
Adaptive HyperLIF uses an input-dependent reset conductance $G(r,J) = (G_0+\beta |J|)r^2$. Full stability and Nengo implementation details are cited to the companion reference paper.

\section{Experimental Design}
Datasets should include the existing event and audio datasets used in the PyTorch benchmark: N-MNIST as the near-static control, DVS-Gesture, SHD, SSC, and any available full-data cache variants.

Architectures:
\begin{itemize}
  \item feedforward SNN;
  \item convolutional SNN;
  \item recurrent SNN;
  \item learnable-delay SNN;
  \item residual/deep convres SNN as a stress test and limitation case.
\end{itemize}

Neuron/readout roster:
\begin{itemize}
  \item LIF, PLIF, ALIF/LSNN, HyperLIF, Adaptive HyperLIF;
  \item HyperLIF variants: $\theta_{\mathrm{event}}=0.95$, origin, dynamic reset;
  \item readouts: membrane, spike count, gate, and hybrid gate/membrane/event readouts.
\end{itemize}

\section{Results Plan}
\subsection{Temporal-Specific Learning Gain}
Report mean test accuracy over seeds. The safe claim is that HyperLIF gains are largest on temporal datasets and collapse toward small gains on N-MNIST.

\subsection{Gate as Decodable State}
Show that LIF gate-only readout is structurally blind while HyperLIF gate readout carries task signal.

\subsection{Accuracy-Spike Efficiency Dial}
Report accuracy versus events per sample. Include Pareto plots and identify operating points, not only best accuracy.

\subsection{Architecture Transfer}
Use the existing architecture-transfer heatmap. Claim transfer across practical architectures; bound the deep convres result honestly if it remains a controlled negative. The current SHD-wide architecture shortlist supports a nuanced baseline claim: HyperLIF leads the recurrent cell, ALIF slightly leads the delay cell, and both are materially above LIF/PLIF under the selected readouts.

\subsection{Learning Rule Robustness}
Present the Julia/APEX gate-eligibility result as an annealed eligibility story: HyperLIF gives a physically meaningful surrogate-temperature path, not a surrogate-free miracle.

\subsection{Cross-Framework Reproducibility}
This section now has calibrated cross-framework evidence. The HDF5 data gate, APEX topology gate, APEX engine import, CUDA availability check, adapter-backed loader smoke, bounded APEX train/eval smoke, and multi-seed tuning pass on byte-identical caches. The safe claim is that tuned APEX HyperLIF/NSA with stronger gate feedback reproduces a small advantage over matched dense LIF; the older default APEX gate setting does not.

\section{Figures and Tables}
\begin{table}[t]
\centering
\caption{Required main-paper artifacts before submission.}
\begin{tabular}{lll}
\toprule
Artifact & Status & Source \\
\midrule
Temporal gain table & Exists, polish & 07 / 08 results \\
Gate readout figure & Exists, polish & 07 figures \\
Efficiency Pareto & Exists, polish & 08 architecture results \\
Architecture heatmap & Exists & 08 paper backbone \\
APEX parity table & Tuned 5-seed APEX comparison passed & Julia HDF5/APEX adapter, tuning grid, and summary CSV \\
Baseline comparison & FF + SHD-wide arch shortlist passed & Final table scope pending \\
\bottomrule
\end{tabular}
\end{table}

\section{Submission Checklist}
\begin{itemize}
  \item Run the expanded PLIF and ALIF/LSNN baseline matrix after smoke validation.
  \item Integrate the tuned APEX GPU validation table; the default-gamma negative result must be framed as calibration evidence, not hidden.
  \item Decide whether the SHD-wide \texttt{rsnn}/\texttt{delay} shortlist is sufficient for reviewer response or whether full-grid PLIF/ALIF tables are required.
  \item Report mean, standard deviation or confidence intervals, paired tests, and failure counts.
  \item Remove unsafe claims about old TNNLS benchmark numbers.
  \item Convert this proposal into full IEEEtran Overleaf source.
\end{itemize}

\section{Conclusion}
This paper should be submitted first after paired statistics, bibliography verification, and advisor-level copyediting are complete. The first IEEEtran source package now exists and builds. Current readiness: approximately 94 percent.

\end{document}
